% TFM-PROTECTED-TIKZ: fig:problema
% Figura fundamental: no eliminar, desactivar, rasterizar ni excluir del PDF.
\providecolor{FeasGreen}{HTML}{2E7D32}
\providecolor{NeedRed}{HTML}{C0392B}
\providecolor{Cardboard}{HTML}{C8954E}
\providecolor{VIUDeep}{HTML}{C63C32}
\providecolor{VIUWhite}{gray}{1}
\providecolor{VIULight}{HTML}{F8FAFC}

% Caja/palé en pseudo-3D: \crate{x}{y}{semilado}.
\providecommand{\crate}[3]{%
  \begin{scope}[shift={(#1,#2)}]
    \fill[Cardboard!35] (#3,#3) -- (#3+0.12,#3+0.13) --
      (#3+0.12,-#3+0.13) -- (#3,-#3) -- cycle;
    \fill[Cardboard!80] (-#3,#3) -- (-#3+0.12,#3+0.13) --
      (#3+0.12,#3+0.13) -- (#3,#3) -- cycle;
    \fill[Cardboard!55] (-#3,-#3) rectangle (#3,#3);
    \draw[Cardboard!95,line width=0.5pt]
      (-#3,-#3) rectangle (#3,#3)
      (#3,#3) -- (#3+0.12,#3+0.13)
      (-#3,#3) -- (-#3+0.12,#3+0.13) --
      (#3+0.12,#3+0.13) -- (#3,#3);
    \draw[Cardboard!95,line width=0.4pt]
      (0,-#3) -- (0,#3) (-#3,0) -- (#3,0);
  \end{scope}}

% Zona de descarga d_k marcada en el suelo.
\providecommand{\dropoff}[3]{%
  \begin{scope}
    \begin{scope}
      \clip[rounded corners=1.5pt]
        ($(#1,#2)+(-0.46,-0.30)$) rectangle ($(#1,#2)+(0.46,0.30)$);
      \fill[VIUDeep!8]
        ($(#1,#2)+(-0.46,-0.30)$) rectangle ($(#1,#2)+(0.46,0.30)$);
      \foreach \i in {-0.5,-0.3,...,0.7}{%
        \draw[VIUDeep!32,line width=0.6pt]
          ($(#1,#2)+(\i-0.45,-0.35)$) -- ($(#1,#2)+(\i+0.10,0.35)$);}
    \end{scope}
    \draw[VIUDeep,line width=0.7pt,rounded corners=1.5pt]
      ($(#1,#2)+(-0.46,-0.30)$) rectangle ($(#1,#2)+(0.46,0.30)$);
    \node[font=\tiny\bfseries,text=VIUDeep,fill=VIUWhite,inner sep=0.5pt]
      at (#1,#2) {#3};
  \end{scope}}

% AMR en vista cenital: \AMR{x}{y}{rumbo en grados}{identificador}.
\providecommand{\AMR}[4]{%
  \pic[rotate=#3] at (#1,#2) {AMRbody};
  \node[circle,fill=VIUWhite,draw=VIUOrange,line width=0.4pt,
        inner sep=0pt,minimum size=0.24cm,font=\tiny\bfseries,text=VIUDark]
        at (#1,#2) {#4};}

\begin{figure}[H]
\centering
\begin{tikzpicture}[x=1cm,y=1cm,>=Stealth,
  AMRbody/.pic={
    \fill[VIUDark!75] (-0.31,-0.14) rectangle (-0.21,0.14);
    \fill[VIUDark!75] ( 0.21,-0.14) rectangle ( 0.31,0.14);
    \draw[draw=VIUOrange,line width=0.6pt,fill=VIUOrange!25,
      rounded corners=1.5pt] (-0.26,-0.20) rectangle (0.26,0.20);
    \fill[VIUDeep] (0.12,-0.11) -- (0.265,0) -- (0.12,0.11) -- cycle;
  },
  pickup/.style={rectangle,rounded corners=2pt,draw=VIUDark!65,
    fill=VIUDark!4,dash pattern=on 3pt off 2pt,minimum width=0.96cm,
    minimum height=0.96cm,inner sep=0pt},
  comm/.style={draw=VIULinkBlue!75,line width=0.7pt,
    dash pattern=on 2.5pt off 1.8pt},
  commrange/.style={draw=VIULinkBlue!55,dotted,line width=0.8pt},
  move/.style={->,draw=VIUOrange!85,line width=0.9pt,
    shorten >=2pt,shorten <=2pt},
  haul/.style={-{Stealth[length=5.5pt,width=4.5pt]},draw=VIUOrange,
    line width=1.5pt},
  route/.style={draw=VIUGray!50,dash pattern=on 3pt off 2.5pt,
    line width=0.7pt},
  phase/.style={rounded corners=3pt,fill=VIUWhite,font=\tiny,
    align=center,inner sep=2.6pt,line width=0.8pt}
]

% Planta y textura del suelo.
\draw[VIUDark!40,line width=1pt,rounded corners=7pt]
  (0,0) rectangle (12.5,8.2);
\begin{scope}
  \clip[rounded corners=7pt] (0,0) rectangle (12.5,8.2);
  \draw[VIUGray!10,line width=0.3pt,step=0.8cm]
    (-0.4,-0.4) grid (12.9,8.6);
  \foreach \sx in {0.35,0.95}{
    \fill[VIUGray!16] (\sx,7.55) rectangle (\sx+0.42,7.95);}
  \foreach \sx in {11.2,11.8}{
    \fill[VIUGray!16] (\sx,7.55) rectangle (\sx+0.42,7.95);}
\end{scope}
\node[font=\tiny,text=VIUGray,anchor=north west] at (0.22,8.06)
  {Almacén automatizado ($20\,\mathrm{m}\times 14\,\mathrm{m}$),
   sin coordinador central};

\node[phase,draw=VIUDark!55,fill=VIUDark!4,text=VIUDark] at (5.75,7.2)
  {base: $|\mathcal C_k|\ge n_k$;\quad general:
   $\sum_{i\in\mathcal C_k}\boldsymbol c_i\succeq\boldsymbol r_k$};

% Disco de comunicación.
\fill[VIULinkBlue!7] (9.25,6.0) circle (1.5cm);
\draw[commrange] (9.25,6.0) circle (1.5cm);
\node[font=\tiny,text=VIULinkBlue!80,anchor=north] at (9.25,4.46)
  {rango de comunicación $R$};

% Zonas de entrega.
\dropoff{1.4}{1.0}{$\boldsymbol d_1$}
\dropoff{11.4}{1.35}{$\boldsymbol d_2$}

% Tareas y zonas de recogida.
\node[pickup] (L1) at (2.0,6.5) {};
\crate{2.0}{6.5}{0.16}
\node[font=\tiny,text=VIUDark,below=1pt of L1]
  {$\boldsymbol l_1\!: k{=}1,\,n_1{=}1$};

\node[pickup] (L2) at (4.6,5.7) {};
\node[font=\tiny,text=VIUGray,opacity=0.9] at (4.6,5.7) {\itshape vacía};
\node[font=\tiny,text=VIUDark,below=1pt of L2]
  {$\boldsymbol l_2\!: k{=}2,\,n_2{=}3$};

\node[pickup] (L3) at (9.9,6.6) {};
\crate{9.9}{6.6}{0.23}
\node[font=\tiny,text=VIUDark,above=1pt of L3]
  {$\boldsymbol l_3\!: k{=}3,\,n_3{=}2$};

\draw[route] (4.6,5.25) -- (7.7,3.95) -- (11.0,1.7);

% Tarea 1: un robot.
\AMR{2.0}{5.15}{90}{1}
\draw[move] (2.05,5.45) -- (2.05,6.05);
\node[phase,draw=FeasGreen,text=FeasGreen,anchor=west] at (2.55,5.2)
  {$n_1{=}1$ \checkmark};

% Tarea 2: coalición en transporte.
\draw[comm] (7.05,3.40) -- (8.35,3.40);
\draw[comm] (7.05,3.40) -- (7.70,4.45);
\draw[comm] (8.35,3.40) -- (7.70,4.45);
\AMR{7.05}{3.40}{-28}{2}
\AMR{8.35}{3.40}{-28}{3}
\AMR{7.70}{4.45}{-28}{4}
\crate{7.70}{3.74}{0.25}
\draw[<->,VIUGray,line width=0.7pt] (7.05,3.02) -- (8.35,3.02)
  node[midway,below,font=\tiny,text=VIUGray,inner sep=1pt] {$d^*$};
\draw[haul] (8.55,3.25) -- (10.55,1.85);
\node[font=\tiny,text=VIUOrange,anchor=west] at (9.45,2.75)
  {transporte};
\node[phase,draw=FeasGreen,text=FeasGreen,anchor=east] at (6.45,2.45)
  {\textbf{transporte}\\ $|\mathcal C_2|{=}3{=}n_2$ \checkmark};

% Tarea 3: reclutamiento.
\draw[comm] (9.35,5.95) -- (8.05,6.35);
\AMR{9.35}{5.95}{55}{5}
\AMR{8.05}{6.35}{-8}{6}
\draw[move] (9.45,6.20) -- (9.65,6.45);
\draw[move] (8.35,6.42) to[out=15,in=205] (9.45,6.78);
\node[phase,draw=NeedRed,text=NeedRed] at (9.25,5.05)
  {\textbf{reclutamiento}\\ $|\mathcal C_3|{=}1<n_3{=}2$};

% Leyenda.
\begin{scope}
\draw[VIUGray!35,rounded corners=4pt,fill=VIULight]
  (0,-1.9) rectangle (12.5,-0.5);
\pic at (0.55,-0.92) {AMRbody};
\node[font=\tiny,anchor=west] at (0.95,-0.92) {AMR (vista cenital)};
\crate{3.55}{-0.92}{0.16}
\node[font=\tiny,anchor=west] at (3.85,-0.92) {carga ($\propto n_k$)};
\node[pickup,minimum width=0.42cm,minimum height=0.42cm] at (6.35,-0.92) {};
\node[font=\tiny,anchor=west] at (6.62,-0.92) {recogida $\boldsymbol l_k$};
\dropoff{9.05}{-0.92}{}
\node[font=\tiny,anchor=west] at (9.6,-0.92) {entrega $\boldsymbol d_k$};
\draw[comm] (0.30,-1.5) -- (0.85,-1.5);
\node[font=\tiny,anchor=west] at (0.92,-1.5) {enlace com.};
\fill[VIULinkBlue!10] (3.1,-1.5) circle (0.16cm);
\draw[commrange] (3.1,-1.5) circle (0.16cm);
\node[font=\tiny,anchor=west] at (3.4,-1.5) {rango $R$};
\draw[haul] (5.55,-1.5) -- (6.2,-1.5);
\node[font=\tiny,anchor=west] at (6.28,-1.5) {transporte};
\node[FeasGreen,font=\scriptsize] at (8.3,-1.5) {\checkmark};
\node[font=\tiny,anchor=west,text=FeasGreen] at (8.48,-1.5)
  {$|\mathcal C_k|{=}n_k$};
\node[NeedRed,font=\scriptsize] at (10.25,-1.5) {$\bullet$};
\node[font=\tiny,anchor=west,text=NeedRed] at (10.42,-1.5)
  {$|\mathcal C_k|{<}n_k$};
\end{scope}

\end{tikzpicture}
\caption[Escenario de asignación distribuida con coaliciones]{Instancia
conceptual del problema de asignación distribuida con formación de coaliciones.
Una flota de seis AMR atiende tres cargas heterogéneas. La tarea 3 permanece en
reclutamiento y la tarea 2 transporta su carga hacia $\boldsymbol d_2$ con
geometría $d^*$. El diagrama representa la arquitectura objetivo local; no
reclasifica como distribuidos los cierres o certificados globales declarados en
los experimentos.}
\label{fig:problema}
\viuownsource
\end{figure}
